\section{Stable recovery from three fixed slices}
\label{sec:rank-two-coercivity}
Frequency-assigned sampling gives a general nonlinear recovery
result. Parallel transport from the two characteristic directions
can separate their contributions in a full orbit snapshot.
This yields local nonlinear stability with a fixed number of
slices near a broad class of inhomogeneous backgrounds,
including the explicit Gowdy-to-Ernst backgrounds of
Section~\ref{sec:rank-two}.

Let $\bar U=(\bar P,\bar Q)$ be a smooth background, use
$t=a+d$, and let $V$ be a linearized wave map with
$V(a)=0$, $D_tV(a)=f e_P$. The smooth orthonormal frame
$e_P=\partial_P$, $e_Q=e^{-\bar P}\partial_Q$ has connection
one-form $\omega=e^{\bar P}d\bar Q$, with
$D_\mu e_P=\omega_\mu e_Q$.
Let $A_\pm(t,\theta)$ be the frame component vector obtained
by transporting $e_P$ along the $\pm$ null rays from
$(a,\theta\mp d)$ to $(t,\theta)$.
For $h=\partial_\theta^{-1}(f-\langle f\rangle)$ of mean zero,
the leading displacement in this frame is
\begin{equation}
 \sqrt{t/a}\,V(t,\theta)
  =\frac12\left[A_-(t,\theta)h(\theta+d)
                    -A_+(t,\theta)h(\theta-d)\right]+R_t f .
 \label{eq:transport-leading}
\end{equation}
For integer $s\ge1$, $R_t$ maps $H^{s-1}$ boundedly to
$H^{s+1}$ on any smooth regular-time interval in question.

To verify the remainder estimate, multiply the covariant
linearization by $\sqrt{t/a}$ to remove $t^{-1}D_t$.
Its principal part is $D_t^2-D_\theta^2$ and its other
terms are a smooth zeroth-order potential given by curvature
and $1/(4t^2)$. Apply this operator to the leading expression
in \eqref{eq:transport-leading}. The terms with $h''$
cancel between equal-speed time and space derivatives.
The terms with $h'$ vanish by the two transport equations
$(D_t\pm D_\theta)A_\pm=0$.
Only smooth coefficients times $h$ remain.
At $a$ the leading expression is zero and its initial
derivative equals $f-\langle f\rangle$ in the $e_P$ direction
up to a smooth coefficient times $h$.
The error therefore has zero initial displacement, initial
velocity bounded in $H^s$ by $\|f\|_{H^{s-1}}$, and source
with that same regularity. The linear energy estimate proves
the claim. The formula is insensitive to the finite-rank
choice of treatment of the mean.


\subsection{A criterion from the relative null transport}
At a chosen terminal time $c=a+d$, let $\alpha_\pm(\theta)$
be the angles of $A_\pm$, defined by the signed integrals
$-\int_{\gamma_\pm}\omega$ on the two null rays.
Write $\rho=(\alpha_- -\alpha_+)/2$, $\beta=\sin\rho$,
and rotate the output frame by the mean angle
$(\alpha_-+\alpha_+)/2$.
The two leading components in \eqref{eq:transport-leading},
up to the common factor $1/2$, then read
\[
 \cos\rho\,\big(H-\mathcal T H\big),\qquad
 \beta\,\big(H+\mathcal T H\big),
 \quad H(\theta)=h(\theta+d),\quad
 \mathcal T H(\theta)=H(\theta-2d).
\]
We use the conditions
\begin{equation}
 \inf_\theta|\cos\rho(\theta)|>0,\qquad
 \inf_{\theta\in\mathbb T}\sum_{j=0}^{J}
           |\beta(\theta-2jd)|^2>0
 \quad\text{for some finite }J .
 \label{eq:rotation-cover}
\end{equation}
For an irrational $d/\pi$ and $\beta\not\equiv0$,
the second condition follows from density of
circle-rotation orbits and a finite subcover of the
translates of an open set where $|\beta|$ is bounded below.

In particular the conditions hold at all sufficiently
short delays after a slice where $Q_\theta(a)$ is
nonzero somewhere. Indeed, expansion of the two ray
integrals about $(a,\theta)$ gives
\begin{equation}
 2\rho=2d\,\omega_\theta(a,\theta)
        +d^2\big(\partial_t\omega_\theta-
                  \partial_\theta\omega_t\big)(a,\theta)
        +O(d^3).
 \label{eq:relative-expansion}
\end{equation}
If $Q_\theta(a)\equiv0$ but $P_\theta(a)Q_t(a)$
is nonzero somewhere, the first term vanishes and the
coefficient of $d^2$ is nonzero on an interval since
$d\omega=e^{P}dP\wedge dQ$. These two cases include any
background of rank two at some point of the initial slice.
In either case a fixed open arc has $|\beta|$ bounded
below by a positive multiple of $d$ or $d^2$ for small
$d$. If the arc is wider than $2d$, the first
$\lceil\pi/d\rceil+1$ translates cover the circle.
The short-delay result therefore holds for rational
as well as irrational placements.

For the explicit Gowdy-to-Ernst construction,
$\omega_t=C e^{-P(w)}w_\theta$ and
$\omega_\theta=C e^{-P(w)}w_t$.
Set $B'(z)=C e^{-P(z)}$ with $B(0)=0$.
The two rotations are $-B(w(c,\theta))$ and
$+B(w(c,\theta))$, since the seed phase vanishes at $a$.
For a small-amplitude seed of a single mode $k$,
a time $c$ with nonzero mode amplitude and
$2kd\notin\pi\mathbb Z$ satisfies both conditions,
with $J=1$.

\begin{theorem}
\label{th:fixed-rank-two}
Consider a smooth Gowdy background and a terminal time
$c$ satisfying \eqref{eq:rotation-cover}. In addition to the initial
$P(a),Q(a),Q_t(a)$ and the metric anchor, read $P,Q$ at $c$
and finitely many low Fourier coefficients of $P$ at a
suitably chosen $a+d_1\in(a,c)$.
There is an $H^s$ neighborhood of the background initial
$P_t$ for each integer $s\ge1$ in which the nonlinear
observation is locally bi-Lipschitz from the missing
velocity in $H^s$ to these observations, using $H^{s+1}$
for the full orbit fields and a fixed finite-dimensional
norm for the early coefficients. The same estimate applies on the
periodic momentum-constraint subspace.
\end{theorem}

\begin{proof}
The displayed principal components control
$H-\mathcal T H$ and $\beta H$ in $L^2$.
If $b_j(\theta)=\beta(\theta-2jd)$, then
\[
 b_j H = \mathcal T^j(\beta H)+b_j(H-\mathcal T^j H).
\]
Telescoping bounds the second term by
$j\|\beta\|_\infty\|H-\mathcal T H\|_{L^2}$.
Summing these estimates and using \eqref{eq:rotation-cover}
controls $\|H\|_{L^2}$.
Apply the same estimate to $\partial_\theta^{s+1}H$.
Commutation through the smooth multiplication coefficients
adds only terms bounded by $\|H\|_{H^s}$.
Add the compact remainder from \eqref{eq:transport-leading}
and the mean mode. For the exact derivative
$\mathcal A_c:f\mapsto (\delta P(c),\delta Q(c))$,
this gives
\begin{equation}
 \|f\|_{H^s}\le C
    \big(\|\mathcal A_c f\|_{H^{s+1}}+\|f\|_{H^{s-1}}\big).
 \label{eq:rank-two-compact}
\end{equation}
The frame factors and the scalar time rescaling
are bounded invertible multipliers on the output space.


The low-frequency reading can be chosen constructively.
Let $C$ be the constant in \eqref{eq:rank-two-compact},
and choose $K$ so that $C/\langle K\rangle<1/4$.
The $H^{s-1}$ norm splits into the low part and a tail
bounded by $\langle K\rangle^{-1}\|f\|_{H^s}$.
For $d_1<\min(d,1/K)$, apply the free-wave Duhamel formula
to the coordinate linearization from $a$ to $a+d_1$.
The map $f\mapsto\Pi_K\delta P(a+d_1)/d_1$
differs from the free diagonal map
$\operatorname{sinc}(k d_1)\Pi_Kf$
by a map bounded from $H^{s-1}$ to $H^{s-1}$ with norm
$O(d_1)$. This follows from the smooth coefficient
energy estimate and $|S_k(r)|\le r$.
The free gains are at least $\sin(1)$.
Choose $d_1$ smaller if needed to absorb
the $O(d_1)\|f\|_{H^{s-1}}$ contribution into the
left-hand side of \eqref{eq:rank-two-compact}.
The joint derivative is bounded below.
Equivalently, compactness in \eqref{eq:rank-two-compact}
shows that the full-terminal kernel is finite dimensional,
and the early reading removes it.

For the nonlinear step, the first-order energy formulation
with coordinate source $\mathcal F$ gives a $C^1$ solution
map on a neighborhood of the smooth background, on the fixed
time interval. One can obtain the derivative by subtracting the
linearized equation from the difference equation and applying
the same energy estimate to the quadratic Taylor remainder.
The derivative is continuous in operator norm by a linear
energy estimate on the difference of two coefficient fields.
Thus on a sufficiently small convex ball the observation
derivative differs from its value at the background by
less than half the lower bound of the joint derivative.
Integrating this operator difference along the segment between
two velocities proves the pairwise lower Lipschitz estimate.
The upper estimate follows from the local forward energy bound.
The affine periodic momentum constraint can be imposed after
the estimate.
\end{proof}

This result applies to full orbit observations at a time
where the two null transports differ.
At a common node of the seed phase, $\rho=0$ everywhere
and the argument supplies only the translation difference.
It does not give a stability estimate at the resonant
two-metric pair. The construction uses two later time
slots, and the earlier slot needs only finitely many
low Fourier coefficients. It leaves open
whether that extra reading can be omitted or replaced
by a particular previously imposed terminal time.

\subsection{When the full terminal orbit snapshot suffices}

\begin{proposition}
\label{prop:two-slice-open}
There is an open class of smooth Cauchy backgrounds,
including solutions with rank-two target differential on
open sets, near which the nonlinear missing-velocity
inverse is locally bi-Lipschitz from one full terminal
orbit snapshot and the initial $P(a),Q(a),Q_t(a)$.
The norms are $H^s$ for velocity and $H^{s+1}$ for the
snapshot, for integer $s\ge1$.
One explicit choice starts at a geodesic-image background
$U=\gamma(w)$ with $\gamma$ a unit-speed nonvertical
geodesic. Suppose $w(a,\theta)$ is nonconstant and its
range lies in a compact segment on which
$|\langle e_P,\dot\gamma\rangle|>0$.
A sufficiently short delay $d=\pi/L$ with $L$ an integer
gives the conclusion.
\end{proposition}

\begin{proof}
On a geodesic-image background the tangential and normal
linearized equations decouple in the frame parallel along
$\gamma$. The curvature terms take values in the normal
direction. Thus the tangent component $X$ obeys
$X_{tt}+t^{-1}X_t-X_{\theta\theta}=0$ with initial
position zero and velocity
$X_t(a)=\langle e_P,\dot\gamma(w(a))\rangle f$.
The full terminal orbit snapshot contains $X(c)$.
The scalar rational-gap bound of
Proposition~\ref{prop:safe-gap} gives, for sufficiently
large $L$,
\[
 \|X_t(a)\|_{H^{s-1}}\le C\|X(c)\|_{H^{s+1}} .
\]
The multiplier $\langle e_P,\dot\gamma(w(a))\rangle$
is smoothly invertible by hypothesis. In addition,
$Q_\theta(a)$ is nonzero somewhere by nonverticality
and nonconstant $w(a)$. The same short delay satisfies
\eqref{eq:rotation-cover}. We may therefore use the
tangent estimate to control the $H^{s-1}$ remainder
in \eqref{eq:rank-two-compact} and obtain a bounded-below
derivative for the single terminal full orbit snapshot.
The $C^1$ argument of Theorem~\ref{th:fixed-rank-two}
turns this into a local nonlinear inverse. Continuity
in the background Cauchy data preserves the strictly
positive lower bound at the same time. A small initial
velocity perturbation perpendicular to the geodesic
preserves the momentum constraint and creates rank two
where $w_\theta(a)\ne0$. This gives genuinely rank-two
members of the open class.
\end{proof}

\paragraph{Necessity of the cover at principal order.}
When $\cos\rho$ is bounded away from zero,
\eqref{eq:rotation-cover} is also necessary for the
one-terminal linearized orbit observation to have an
energy lower bound, before any lower-frequency
obstruction is considered. Indeed, if the cover fails,
compactness gives a point whose full forward rotation
orbit lies in the zero set of $\beta$.
For irrational rotation this forces $\beta\equiv0$,
and the translation difference has arbitrarily small
Fourier gains by rational approximation.
For rational rotation the zero orbit is finite.
Choose smooth, rotation-invariant envelopes supported in
shrinking neighborhoods of that orbit and multiply them
by arbitrarily high Fourier carriers invariant under the
same rotation. The principal translation difference
vanishes and the multiplication by $\beta$ becomes
arbitrarily small. Taking the carrier large for each
fixed envelope makes its contribution from the
lower-order remainder tend to zero in the normalized
energy scale. These sequences contradict the required
lower bound. This argument concerns the full orbit
snapshot without an additional terminal $\lambda$ reading.
